% Figaro Paper H: The Coercion Variable
% On the Boundary Between What the Sovereign Must Do
% and What the Bond Can Do
% (Political-philosophy treatment of coercion as a substrate variable;
% companion to F1/F2 on civil-law subjecthood and to B3 on subordination.)
% Preprint, April 2026

\documentclass[11pt,a4paper]{article}

\usepackage[utf8]{inputenc}
\usepackage[T1]{fontenc}
\usepackage{amsmath,amssymb,amsthm}
\usepackage{mathtools}
\usepackage{booktabs}
\usepackage{graphicx}
\usepackage{hyperref}
\usepackage{cleveref}
\usepackage{enumitem}
\usepackage{xcolor}
\usepackage[margin=1in]{geometry}
\usepackage{natbib}
\bibliographystyle{plainnat}

\theoremstyle{remark}
\newtheorem{remark}{Remark}[section]

\newcommand{\figaro}{\textsc{Figaro}}
\newcommand{\fig}{\textsc{Fig}}

\title{%
  The Coercion Variable:\\
  On the Boundary Between What the Sovereign Must Do\\
  and What the Bond Can Do%
}

\author{
  Alessandro Daliana\thanks{Corresponding author. \texttt{adaliana@figaro.org}}
}

\date{April 2026}

\begin{document}

\maketitle

% ════════════════════════════════════════════════════════════════════
\begin{abstract}
The dominant tradition of legal philosophy from \citet{austin1832}
and \citet{hobbes1651} through \citet{weber1919} and
\citet{hart1961} has treated law as a system of texts paired with
a coercive enforcement apparatus. The texts specify obligations;
the apparatus --- police, courts, sheriffs, ultimately the
sovereign's monopoly on legitimate violence --- ensures the
obligations are met or punishes their breach. The exact
articulation of the relationship between text and coercion varies
across the tradition (command theory, secondary rules, legitimate
violence, and so on), but the structural pairing is shared.

This paper observes that the bonded settlement primitive examined
in the corpus's companion mechanism, institutional-economics, and
political-economy papers performs the enforcement function of
law in bilateral commerce \emph{without} a coercive apparatus in
the tradition's sense. There is no third party who applies force
upon breach. There is no court that compels performance. There is
no sovereign whose monopoly on violence stands behind the
arrangement. The enforcement is performed by collateral that the
parties themselves lock at the moment of consent, and that the
protocol releases or forfeits according to the parties' subsequent
behavior. The mechanism is, in the precise language of
\citet{elster1979ulysses}, a \emph{pre-commitment} structure:
self-binding at the moment of consent, with no external party in
the position to apply or withhold force after the fact.

Pre-commitment via bond is qualitatively different from
retributive coercion via the sovereign. The difference is not a
matter of degree --- a softer or kinder enforcement --- but of
mechanism: the relationship between obligation and enforcement
runs through the parties' own resources and consent, not through
an external authority's standing capacity to compel. Where the
mechanism applies, it substitutes for the coercive apparatus.
Where it does not apply --- against criminal harm, against
externalities affecting non-parties, against the provision of
public goods, against status and rights adjudication --- the
coercive apparatus continues to do work that the bonded mechanism
is structurally unable to do.

The substantive claim of this paper is therefore not that the
sovereign is unnecessary or that the state should be reduced. It
is that the political-philosophical question of legitimate
coercion --- which thinkers from Hobbes to Schmitt have addressed
as if it ranged over the entirety of social cooperation ---
becomes a bounded question once a substantial part of bilateral
commerce can be coordinated by pre-commitment rather than by
sovereign-backed coercion. The state retains its load-bearing
functions; the question of how it should exercise them retains
its political-philosophical urgency; what changes is the
\emph{domain} over which the question must be answered. This
relocation is the political-philosophical content of the
substrate change, and it is the subject of this paper.
\end{abstract}

\medskip
\noindent\textbf{Keywords:}
coercion, sovereignty, monopoly on violence, pre-commitment,
bonded enforcement, political philosophy, Hobbes, Weber, Hart,
boundary of state authority

% ════════════════════════════════════════════════════════════════════
\section{Introduction}
\label{sec:intro}

This paper sits at the intersection of political philosophy and
the analysis of a particular coordination primitive. Its
starting point is an observation that emerged in conversation
about how to characterize the bonded settlement primitive
relative to the dominant Western tradition on the nature of law.
The observation is sharp enough to warrant explicit treatment.

The dominant tradition reads law as the pairing of texts (the
obligations) with a coercive apparatus (the enforcement). The
coercive apparatus performs two functions: it compels performance
of obligation when a party would not otherwise perform, and it
punishes breach as a deterrent against future breach by the same
party or by others. The apparatus is grounded in the sovereign's
monopoly on legitimate violence within a territory
(\citealt{weber1919}). When the obligation is honored, the
apparatus does no observable work; when the obligation is
breached, the apparatus is invoked. The apparatus's standing
capacity is what makes the texts effective even in cases where
it is not invoked: the threat of enforcement does the work even
when the actual enforcement is not required.

The bonded settlement primitive examined in the corpus's
companion mechanism paper performs the enforcement function for
bilateral commercial obligations \emph{without} a coercive
apparatus in the tradition's sense. The bond is locked at the
moment of consent, by the parties themselves, from their own
resources. The protocol's resolution mechanic releases the bond
to the cooperative party and forfeits it from the defecting
party, automatically and without discretion. There is no
sovereign behind the arrangement, no court to which a party can
appeal for force, no agent in the position to apply or withhold
enforcement after the fact. The enforcement is internal to the
arrangement; the only external requirement is the public
infrastructure on which the protocol runs.

This is not a marginal substitution. In bilateral commercial
exchange, the coercive apparatus has historically been the
guarantor of last resort against which all other enforcement
mechanisms (reputation, repeated interaction, social pressure)
operate. Removing the requirement for that guarantor changes the
structural role those secondary mechanisms play. More
fundamentally, it changes what kind of question the philosophical
analysis of coercion is. We argue that this change is the
substantive political-philosophical content of the substrate
change.

\paragraph{The thesis.}
The thesis is structural and bounded. (1) Coercion is a substrate
variable: in the dominant tradition, it is structurally
load-bearing in the enforcement of bilateral commercial
obligations. (2) The bonded primitive removes coercion as a
substrate variable in this domain by substituting pre-commitment
for retribution. (3) The substitution does not extend to domains
in which the coercion was performing functions the bond cannot
perform: criminal harm, third-party externality, public-goods
provision, status, rights adjudication. (4) Where the
substitution applies, the political-philosophical question of
how to legitimate coercive power becomes irrelevant rather than
corrected --- the coercion is no longer being exercised, so the
question of its legitimacy does not arise. (5) Where the
substitution does not apply, the political-philosophical question
retains its full urgency, now operating over a bounded rather
than an unbounded domain.

\paragraph{What this paper is and is not.}
This paper is a short essay in the political philosophy of
coercion. It does not extend the underlying mechanism, the
institutional argument, or the political-economy diagnostic
developed in the companion papers. It takes those contributions
as given and asks what their joint implication is for the
philosophical analysis of coercion that has organized political
thought from Hobbes to the present.

The paper takes no position on whether the state should be
larger, smaller, more interventionist, less interventionist,
democratic, republican, or any other configuration. Those
questions are real, and they survive the substrate change in the
domain where coercion remains load-bearing. The paper's claim is
that the substrate change shrinks that domain by removing
bilateral commerce from it, and that this shrinkage is the
political-philosophical content of the change. The shrinkage is
neither a libertarian victory nor a statist defeat --- it is a
relocation of where the question of legitimate coercion must be
answered.

\paragraph{Position relative to the corpus.}
This paper is part of the law-and-civil-authority cluster (with
the companion labor-law paper on subordination as legal axis and
the companion displaced-populations paper on civil-status
preconditions). It addresses what those papers presuppose: the
coercive apparatus that historically stands behind both labor
contracts and civil status, and what the substrate change implies
for that apparatus when its commercial-enforcement function is
substituted by bond. The political-economy trilogy in the B-series
addresses adjacent questions (institutional-economic mechanism in
B1, hegemonic visibility in B2, conditionality of the left/right
partition in B3). The present paper addresses the
political-philosophy question that the B-series brackets: what is
removed when coercion is removed, and what continues to require
it.

% ════════════════════════════════════════════════════════════════════
\section{Coercion as the Third Substrate Variable}
\label{sec:third-variable}

The companion partition-conditionality paper identified two
substrate variables --- subordination and coordination rent ---
on which the disagreement between the cooperative-Marxian and
classical-liberal traditions had been organized. The bonded
primitive removes both as load-bearing variables in bilateral
commerce; the disagreement does not vanish but relocates from the
substrate to the graph. The present paper identifies coercion as
a third such variable, structurally adjacent to the first two and
analyzed by a different tradition.

\paragraph{Coercion in the dominant tradition.}
\citet{austin1832} formulated the command theory of law: laws are
commands of a sovereign, backed by sanctions for non-compliance.
The pairing of command with sanction is constitutive of law on
this view; an obligation without an enforcement mechanism is, in
Austin's terms, a moral injunction rather than a legal one.
\citet{hart1961} refined the analysis substantially, distinguishing
primary rules of obligation from secondary rules of recognition,
change, and adjudication, and arguing that not all law operates
by sanction. But Hart retained coercive enforcement as one
component of legal systems and acknowledged that, in practice,
the enforceability of legal rules depends on the standing
capacity of officials to apply force when required.
\citet{hobbes1651} provided the upstream
political-philosophical justification for the standing capacity:
covenants without the sword are but words, and of no strength to
secure a man at all.

\paragraph{Weber's distinction between \emph{Macht} and \emph{Herrschaft}.}
\citet{weber1922economy} sharpens the picture in a way the
present argument relies on heavily. \emph{Macht} is the bare
capacity to bring about an outcome despite resistance --- power
in the most general sense. \emph{Herrschaft}, often translated
as ``domination'' or ``authority,'' is a narrower category:
\emph{Macht} that is exercised through a stable relationship
of command and obedience, in which one party has a recognized
right to issue directives and another party a recognized
obligation to comply. Weber's three pure types of legitimate
\emph{Herrschaft} --- traditional, charismatic, and rational-legal
\citep{weber1919} --- describe how command-relationships
are stabilized as authority. The Hobbesian sword Austin and Hart
later analyze is \emph{Herrschaft}-shaped: it presupposes the
sovereign's recognized authority to exercise the coercive
capacity, not just the bare capacity itself.

The bonded primitive does coercion-like work --- the breaching
party's bond is forfeit despite their preference --- but it does
that work in the \emph{Macht} register without reaching for
\emph{Herrschaft}. There is no party whose recognized authority
the forfeiture depends on; the bond is released or forfeit by
public algorithm whose execution the parties have themselves
authorized at the moment of consent. The primitive supplies the
forcing function (\emph{Macht}) without supplying or requiring
the authority structure (\emph{Herrschaft}) that the dominant
tradition has assumed must accompany it. This is what makes
coercion a variable: separating \emph{Macht} from
\emph{Herrschaft} has been historically difficult precisely
because the available enforcement architectures bundled them
together; the bonded primitive is one of the first arrangements
in which the bundle can be unbundled at the level of bilateral
commerce.

The functional locus of coercion across these accounts is the
same: it is what makes contractual obligation effective in cases
where a party would not otherwise perform. Without it, parties
who can defect at lower cost than they can perform will defect,
and the system of obligation collapses into a set of pious
recommendations. With it, performance becomes the rational
strategy because the cost of breach exceeds the cost of
compliance.

\paragraph{Coercion in commercial law specifically.}
Commercial law inherits the dominant tradition's structure.
Contract enforcement runs through the courts; specific
performance and damages awards rely on the state's capacity to
compel transfer of property or impose financial penalties; the
standing threat of these remedies is what makes commercial
contracts effective even in the cases where the remedies are
not invoked. \citet{williamson1985}'s transaction-cost analysis
treats the cost of using courts as one of the components of
market enforcement that hierarchy (the firm) is sometimes able
to economize on; the analysis assumes the courts are available as
the residual enforcement mechanism.

The commercial-law apparatus is not equally available across all
cases. Cross-border transactions involve choice-of-law and
choice-of-forum complexities that often render litigation
prohibitively expensive. Transactions between parties in
jurisdictions whose courts do not recognize each other's
judgments lack effective remedy. Transactions between
informally-constituted parties, parties without legal personality
in the relevant jurisdiction, or parties whose identity is hard
to establish, are similarly under-served. These under-served
cases have historically been managed by intermediaries
(merchant houses, factors, escrow agents) or by extralegal
arrangements (reputation, repeated interaction). Each of these
secondary mechanisms operates in the shadow of the residual
state-coercion option, even when the option is not directly
available --- because the residual option's existence has
shaped the social conventions on which the secondary mechanisms
rely.

\paragraph{Coercion as a variable.}
Treating coercion as a \emph{variable} that can in principle be
set to zero is, like treating subordination as a variable, a move
that the dominant tradition has not historically been able to
make. The reason is that no coordination architecture has been
available in which bilateral commercial obligations could be
enforced at scale without a residual coercive backstop. The bond
posted by the parties under prior arrangements was either (a)
held by an intermediary who would release it to the cooperative
party --- in which case the intermediary's discretion was the
coercive mechanism, simply relocated; or (b) held by a court ---
in which case the court's standing capacity was the coercive
mechanism, simply made explicit; or (c) absent --- in which case
the obligation was un-enforced and the parties bore the
performance risk. The dominant tradition's analysis was correct
on its own terms: under the available coordination architectures,
some form of coercive backstop was structurally necessary.

The bonded settlement primitive is a coordination architecture
in which the bond is held neither by an intermediary nor by a
court, and in which the release is performed by a public
algorithm with no discretion. This is the architectural change
that makes coercion a variable. The variable can now in
principle be set to zero in the bilateral-commerce domain, and
the question of what happens to the dominant tradition's
analysis when that occurs is the question the present paper
addresses.

\begin{remark}[On reading coercion as a variable]
The framing tracks the framing of subordination in the
companion partition-conditionality paper. Both moves rely on
the fact that political philosophy and political economy have
historically treated certain features of social cooperation as
structural constants because no architecture admitted their
removal. When an architecture does admit their removal, the
features become variables, and the analyses that took them as
constants relocate accordingly. The relocations are different
in each case because the analyses were different, but the
underlying move (constant becomes variable; analysis relocates)
is the same.
\end{remark}

% ════════════════════════════════════════════════════════════════════
\section{What the Primitive Does to Coercion}
\label{sec:what-primitive-does}

We now turn from the structural framing to the specific mechanism
by which the bonded primitive substitutes for the coercive
apparatus in bilateral commerce. The mechanism is established in
the companion mechanism paper; we summarize only what is necessary
for the present argument.

The substitution rests on two compositions in the kernel.
\emph{Asymmetric bonding} (Mechanism~1) supplies the credible
threat of loss that makes performance the rational strategy ---
the substitute for the sovereign's coercive sword in the bilateral
case. \emph{Buyer dominance with atomic resolution} (Mechanism~2)
supplies the non-discretion that makes the substitute architectural
rather than dependent on a trusted enforcer: only the buyer can
trigger resolution, and resolution settles every active order
together or not at all, with no agent --- arbitrator, governance
body, oracle, jury --- in a position to defer, partial-resolve, or
selectively enforce. The structural difference between Hobbesian
sovereignty and the bonded primitive is not just that the parties
post their own swords; it is that no third party stands in the
position from which the sovereign formerly acted.

\paragraph{Pre-commitment instead of retribution.}
The bond is locked at the moment of consent, from the parties'
own resources, by their own signatures. The amount is fixed by
the protocol's bonding ratios (twice the contract price for the
buyer, twice the cumulative seller value for the seller), and is
verifiable by both parties before the commitment is signed. The
protocol's resolution mechanic releases the bond to the
cooperative party and forfeits it from the defecting party,
automatically, with no discretion exercised by any human agent
or external authority --- the second mechanism above is what
guarantees that no such discretion exists. The mechanism is what
\citet{elster1979ulysses} analyzed under the heading of
\emph{pre-commitment}: a self-binding move at the moment of
consent that substitutes for ongoing enforcement during
performance.

Pre-commitment differs from retributive coercion in three
structural ways. First, it operates at the moment of consent
rather than at the moment of breach: the binding occurs when the
parties agree, not when one of them tries to deviate. Second,
the resource that performs the binding is the party's own,
posted voluntarily, rather than the sovereign's standing
capacity to apply force. Third, no agent is in the position to
apply or withhold the binding after the fact: the protocol
executes the release or forfeiture mechanically, without
discretion, and without the structural capacity to be lobbied,
captured, or selectively enforced.

\paragraph{Hobbes's covenant with self-applied sword.}
Hobbes's famous remark that ``covenants, without the sword, are
but words'' identifies the structural problem the dominant
tradition has spent four centuries solving by erecting and
legitimating sovereign coercive apparatuses. The bonded
primitive offers a different solution: the covenant arrives with
its own sword, applied automatically by the protocol upon breach,
with the sword being the bond the breaching party has itself
posted. The Hobbesian framing is not refuted; it is satisfied by
a different mechanism. The covenant has its sword; the sword is
not the sovereign's; the parties supply it themselves at the
moment of agreement.

This is a literal reading, not a metaphorical one. The bond is
the sword in the precise sense Hobbes intended: a credible threat
of loss whose existence makes performance the rational strategy.
What changes is who holds and who applies the sword. Under the
sovereign's monopoly, the sovereign holds and applies; the
parties' obligations are effective because the sovereign stands
behind them. Under bonded pre-commitment, the parties themselves
hold (in escrow with the protocol) and the protocol applies; the
parties' obligations are effective because they have themselves
made the breach more costly than the performance.

\paragraph{What this is not: a softer kind of coercion.}
A reader might object that pre-commitment is just coercion with
a friendlier name --- the bond is forfeited against the breaching
party's preferences, after all, and that is coercive in the
ordinary-language sense. The objection elides a structural
difference that the political-philosophical literature has long
recognized. \citet{nozick1969coercion} distinguished proposals
from threats by reference to the baseline against which
compliance is evaluated: a threat alters the baseline in a way
that compels compliance, while a proposal offers an alternative
to the baseline that the recipient may accept or refuse. The
bond posted at consent is a proposal in Nozick's sense: at the
moment of bonding, the party's baseline is unchanged; the bond
arrives only with the party's signature, and the consequences of
forfeiture are baked into the offer. This is the ordinary
structure of voluntary commerce. It is categorically different
from the standing capacity of a sovereign to apply force against
a party who has not consented.

\begin{remark}[On the limit of the analogy]
The pre-commitment / retribution distinction is structural,
not ethical. We are not claiming that pre-commitment is morally
superior to retributive coercion --- that is a normative
question on which the substrate takes no position. We are
claiming the two are different mechanisms of enforcement, with
different agents, different timing, different relations to
consent, and therefore different political-philosophical
analyses. The substitution of pre-commitment for retribution in
a bounded domain therefore changes the political-philosophical
analysis applicable in that domain. This is descriptive, not
prescriptive.
\end{remark}

% ════════════════════════════════════════════════════════════════════
\section{Distance and the Menu of Coercive Alternatives}
\label{sec:distance}

A natural question is when, in practice, the bonded mechanism is
the operative enforcement mechanism and when it is one of several
coexisting mechanisms. The answer is governed by what we will
call \emph{social distance} between the parties, by which we
mean their distance from one another along three dimensions:
physical-geographic, jurisdictional-legal, and
identity-relational.

\paragraph{Physical-geographic distance.}
Parties in physical proximity have access to enforcement
mechanisms that do not require either bond or sovereign:
in-person renegotiation, social pressure, recourse to local
intermediaries, and at the limit, direct physical compulsion.
The bonded mechanism is available to them but is not the only
option. Parties at physical distance lose access to most of these
alternatives. Direct compulsion becomes infeasible; in-person
renegotiation cannot occur; social pressure operates only insofar
as the parties share a community that traverses the distance.

\paragraph{Jurisdictional-legal distance.}
Parties under a shared sovereign jurisdiction have access to
that sovereign's coercive apparatus: courts, contract
enforcement, specific performance, damages, criminal sanction
for fraud. Parties under distinct sovereigns may or may not have
access, depending on the existence of mutual-recognition
treaties, the cost of cross-border litigation, and the
political-economic relationship between the sovereigns. Parties
under sovereigns that do not recognize each other's judgments,
or under no sovereign at all, lack the apparatus entirely.

\paragraph{Identity-relational distance.}
Parties who know each other, who have transacted before, or who
share membership in a small community have access to reputational
enforcement: the cost of breach includes the loss of future
transactions and standing within the community. Parties who are
strangers to each other, who have no prior interaction, and who
share no community lack this. The reputational mechanism does
not extend to first encounters between unknown parties.

\paragraph{The menu argument.}
At low social distance along all three dimensions, parties have
a rich menu of enforcement mechanisms: physical, sovereign-legal,
reputational. The bonded mechanism is one option among several
and is not necessarily the dominant one. At high social distance
along all three dimensions, the alternatives drop away: physical
compulsion is infeasible, sovereign apparatus is unavailable or
prohibitively expensive, and reputation cannot operate between
strangers. The bonded mechanism is then the only enforcement
mechanism available, and its existence becomes the difference
between transaction and no transaction.

This is not an empirical claim that the bonded mechanism will be
adopted at high social distance and ignored at low social
distance --- adoption is a separate question --- but a structural
observation about where the mechanism is doing decisive work and
where it is one option among others. The decisive cases are the
ones in which prior coordination architectures had no answer:
cross-border commerce between parties with no shared sovereign,
exchange between parties whose identities cannot be reliably
established, transactions between artificial agents acting under
distributed organizational arrangements, commerce by
populations who lack civil-legal personality in any jurisdiction
that the counterparty can sue in.

\begin{remark}[On the physical-proximity intuition]
A common intuition holds that the bonded mechanism's incentives
are stronger at distance and weaker in proximity. The intuition
is correct in its direction but imprecise in its mechanism. The
bonded mechanism's incentives do not change with distance ---
the bonding ratios are constant. What changes with distance is
the menu of competing mechanisms. At distance, the menu shrinks
to the bonded mechanism alone, making it decisive. In proximity,
the menu expands to include alternatives, of which the bonded
mechanism is one. The mechanism is not stronger at distance; it
is exclusive at distance and one-of-several in proximity.
\end{remark}

% ════════════════════════════════════════════════════════════════════
\section{The Boundary: What Shrinks, What Stays}
\label{sec:boundary}

The bonded mechanism's domain of substitution is bounded. Within
the boundary, it substitutes for sovereign-backed coercion; the
coercion is no longer being applied because the obligation is
self-enforcing. Outside the boundary, the bonded mechanism is
structurally unable to substitute, and sovereign-backed coercion
continues to perform functions that the bond cannot. This
section identifies what is on each side of the boundary, and why.

\paragraph{Inside the boundary: priceable bilateral commercial obligations.}
The bonded mechanism substitutes for coercion in obligations that
are (a) priceable in a unit the protocol denominates bonds in,
(b) bilateral or composable into bilateral edges via the
mechanism's progressive collateralization, (c) voluntary at the
moment of bonding, and (d) capable of being specified in
schemata that the parties accept at consent. The bulk of
contemporary commercial activity falls under these conditions:
sale of goods, provision of services, delivery, escrow, payment,
attestation against measurable standards. For these, the bonded
mechanism removes the requirement for an external coercive
backstop.

The substitution has cascading effects on the regulatory
apparatus that historically performed coordination functions
adjacent to direct enforcement. Antitrust regimes assume the
existence of intermediaries with market power to constrain;
where intermediation is structurally unnecessary, the
intermediary-targeting regime has less to operate on.
Consumer-protection regimes assume an information asymmetry
between merchant and consumer that schema-validated attestations
substantially reduce. Cross-border commercial-dispute regimes
assume the necessity of choosing between competing sovereign
coercive apparatuses; where the obligation is self-enforcing,
the choice is moot. The shrinkage is real but is itself
bounded: each of these regimes addresses other concerns
(systemic risk, fraud, market manipulation) that survive the
substrate change.

\paragraph{Outside the boundary: non-priceable, non-bilateral, non-voluntary harms.}
The bonded mechanism cannot substitute for sovereign-backed
coercion in cases where the obligation is non-priceable
(violence against persons cannot be made commensurable with a
sum of money in a way that captures what the violation
involves), non-bilateral (third parties harmed by an externality
are not parties to the bond), or non-voluntary (the harming
party did not consent to a bond at the moment of harm). These
include criminal harm against persons, large-scale environmental
externality, public-goods provision under free-rider conditions,
status questions (citizenship, legal personality, identity),
fundamental rights adjudication, and the standing capacity to
defend the territory against external coercion.

For each of these, the sovereign's coercive apparatus continues
to do work that the bonded mechanism is structurally unable to
do. The reason is not that the bonded mechanism is inefficient
in these domains but that the conditions for its operation are
not present: there is no consent at the moment of harm, no
bilateral counterparty to bond against, no priceable
specification of the obligation. Substituting bond for sovereign
in these domains would require the substitution to itself be
sovereign-backed, which is not a substitution.

\paragraph{The boundary itself.}
The boundary between the two domains is not a moral or political
choice but a structural one, given by the conditions under which
bonded pre-commitment is operable. The boundary will shift over
time as schema design improves and as bondable composition
extends to obligations that are not currently priceable in
practice. But the structural limits --- non-consent at the moment
of harm, externality on non-parties, non-priceability of certain
violations --- are categorical and do not yield to engineering
improvement. They mark the part of the cooperative landscape
where sovereign-backed coercion remains load-bearing.

\paragraph{What this implies for the apparatus.}
The state's coercive apparatus is not made unnecessary by the
substrate change. It is bounded, in the precise sense that some
of its prior functions (commercial enforcement of bilateral
obligations) are no longer required for the obligations to be
honored, while other functions (criminal, externality, public
goods, status, rights, defense) continue to require it. The
shrinkage is in the part of the apparatus's work that was
addressing the coordination problem; the part addressing the
not-coordination problems remains.

\begin{remark}[On contested boundary cases]
Some categories sit on the boundary and admit dispute about
which side they fall on. Securities regulation is one example:
some of its function is intermediary-targeting (and shrinks
with the substrate change), some is externality-pricing
(systemic risk to non-parties, which does not shrink), some is
status-based (defining qualified investor categories, which
does not shrink). Disentangling these is empirical work the
present paper does not undertake. We note only that the
boundary is real, that it bounds the political-philosophical
question, and that contested cases at the boundary are
adjudicable on structural grounds rather than political ones.
\end{remark}

% ════════════════════════════════════════════════════════════════════
\section{The Democratic-Republican Apparatus: Bounded Domain, Not Reformed}
\label{sec:apparatus}

We now turn to what the substrate change implies for the
democratic-republican apparatus through which contemporary states
exercise their coercive monopoly. The implication is narrower
than may be tempting to claim, and the framing matters.

\paragraph{The levers of power and their abuses.}
The standard analysis identifies four broad levers of power in
democratic-republican constitutional systems: legislation
(rule-making), executive enforcement, judicial adjudication, and
the franchise (the citizenry's authorization of the previous
three). Each is subject to characteristic abuses: regulatory
capture in legislation (\citealt{stigler1971economic}), executive
overreach and emergency powers
(\citealt{schmitt1922politicaltheology}), judicial drift,
gerrymandering and voter suppression in the franchise, and
across-lever phenomena like agency capture and lobbying
(\citealt{olson1965logic}). These abuses are real, are
extensively documented, and have generated centuries of
political-philosophical and constitutional reflection on how
to constrain them.

\paragraph{What the substrate change does not do.}
The substrate change does not correct the abuses listed above.
Regulatory capture continues to operate in the domains where
regulation continues to operate; executive overreach continues to
be possible in the domains where the executive continues to
exercise authority; judicial drift continues in the domains where
courts continue to adjudicate. The substrate change does not
provide a constitutional remedy for any of these. It is not a
proposal for democratic reform.

\paragraph{What the substrate change does do.}
The substrate change shrinks the \emph{domain} over which the
levers of power operate. Where bilateral commercial obligations
become self-enforcing, the legislative apparatus that has
historically generated rules for their regulation has less to
generate rules about. Where commercial disputes resolve
atomically by protocol, the judicial apparatus that has
historically adjudicated them has less to adjudicate. Where
coordination occurs by bond, the executive apparatus that has
historically enforced commercial obligations has less to enforce.
The abuses that operate on these levers are not corrected; the
levers' \emph{reach} is narrowed in the part of social
cooperation that becomes self-enforcing.

This is subtraction, not reform. The political-philosophical
work of constraining the levers in their remaining domains
continues. The work is, however, more bounded than it was: the
domain over which legitimate coercion must be authorized,
exercised, and constrained is smaller, because a substantial
part of what historically required authorization, exercise, and
constraint is now operating without coercion in the relevant
sense.

\paragraph{Pettit's freedom-as-non-domination, and where it lands here.}
\citet{pettit1997republicanism} reframes liberty in a way the
present argument can use directly. The classical liberal reading
takes liberty to be the absence of \emph{actual} interference; an
agent is free to the extent that no one is currently obstructing
their actions. Pettit's republican reading takes liberty to be the
absence of \emph{capacity} for arbitrary interference; an agent is
free to the extent that no one is in a position to interfere
arbitrarily, even if no one is currently doing so. The republican
reading is a stronger condition: a slave with a kind master is not
free in the republican sense, because the master retains the
capacity to interfere arbitrarily even when not exercising it.

The bonded primitive's non-discretion is a freedom-as-non-domination
property. At the moment of resolution, no third party is in a
position to interfere arbitrarily with the outcome --- not the
buyer (whose dominance is bounded by the conservation law and the
atomic-resolution constraint), not the seller, not a manager, not
a foundation, not a court (the court can adjudicate on the
evidence but cannot reach into the bonded escrow). The
participants in a bonded process enjoy republican-style freedom
within the process's scope: not because no one chooses to
interfere, but because the architecture has eliminated the
position from which arbitrary interference could be exercised.
The connection should not be over-stated --- republican freedom
is a property of political order broadly, and the bonded
primitive is a settlement architecture in a specific domain ---
but the structural alignment is direct, and the primitive supplies
within its domain what republican theory has historically required
political institutions to supply across domains.

\paragraph{What about democracy itself.}
A democratic-republican order has multiple justifications, only
some of which depend on the order's coordinative function in
commerce. The order's role in legitimating coercion (Locke's
consent of the governed), in protecting rights against arbitrary
authority (Madisonian checks), in providing for collective
self-determination (the republican tradition from Cicero through
\citet{pettit1997republicanism, pettit2012on}), and in adjudicating
among competing visions of the good
(the deliberative tradition from Mill through Habermas), are
each independent of the order's role in regulating commerce.
None of these is affected by the substrate change. The order's
coordinative function in commerce, however, is one of its
historically major functions, and that function shrinks. A
democratic-republican order is not made unnecessary by the
substrate change; it is freed, in the bounded domain of
self-enforcing commerce, from one of the responsibilities that
has historically occupied it.

\begin{remark}[On the Network State and adjacent proposals]
A separate question, on which we take no position here, is
whether the substrate change implies that some functions of the
democratic-republican order can or should be replaced by
network-organized arrangements (\citealt{srinivasan2022network}).
That question requires its own treatment and is the subject of
a forthcoming companion paper. The present paper notes only
that the substrate change does not by itself entail any
particular network-organizational form, and that conflating the
substrate change with a specific political proposal is the
category error this paper has been at pains to avoid.
\end{remark}

% ════════════════════════════════════════════════════════════════════
\section{What This Paper Is Not}
\label{sec:not}

The argument is liable to several misreadings; we address them
explicitly.

\paragraph{Not anarchism.}
The paper does not argue that the state is unnecessary, that the
sovereign's monopoly on legitimate violence should be dissolved,
or that the coercive apparatus should be abolished. The state's
load-bearing functions outside the bilateral-commerce domain are
preserved by the analysis. A reader who concludes that the paper
endorses anarchism has read shrinkage in one domain as
abolition across all domains.

\paragraph{Not libertarianism.}
The paper does not argue that the state \emph{should} be
smaller, that voluntary exchange is morally superior to
authoritative coordination, or that the night-watchman state of
\citet{nozick1974anarchy} is the right institutional
configuration. The paper makes no normative claim about whether
the shrinkage is good or bad, only the structural claim that it
occurs in a bounded domain. A reader who concludes that the paper
endorses libertarianism has read structural observation as
normative recommendation.

It is worth distinguishing the present argument from the
adjacent libertarian-anarchist program. \citet{nozick1974anarchy},
\citet{friedman1973machinery}, and the broader anarcho-capitalist
tradition argue that the state's enforcement function can in
principle be replaced by competing private protection agencies,
and that the resulting arrangement would be morally and practically
superior to the state. The structural claim in those programs is
\emph{unbounded substitution}: every coercive function the state
performs can in principle be performed by some private
arrangement, including criminal law, defense, and adjudication of
disputes between strangers. Our claim is the structurally weaker
\emph{bounded substitution}: a specific class of bilateral
commercial obligations becomes self-enforcing under the bonded
primitive, and \emph{outside that class} the conventional
sovereignty apparatus remains in place and is not analyzed here.
The substitution is bounded because the primitive's non-discretion
is purchased by accepting the constraint that the substituted
function be specifiable in advance; the criminal-law function and
the broader adjudicative functions of the state are not so
specifiable, and the bonded primitive does not pretend to
substitute for them. A reader sympathetic to the
libertarian-anarchist program should note that the present
argument supplies neither the unbounded substitution they may
hope for nor a refutation of the desirability of such a program;
it supplies a bounded substitution in one domain and is silent
on the others.

\paragraph{Not statism.}
Symmetrically, the paper does not argue that the state's
remaining functions are well-performed, well-justified, or that
the apparatus that performs them is in good order. The paper is
silent on these questions. The fact that the bonded mechanism
does not substitute for the state's criminal-law function does
not entail that the state's criminal-law function is being
performed correctly. That is a separate inquiry on which the
substrate takes no position.

\paragraph{Hobbes and Schmitt: how the substitution does and does not engage them.}
Hobbes and Schmitt are the two thinkers whose challenges to the
present argument are most often raised, and they require
different responses.

Hobbes's structural observation, that covenants without swords
are but words, is satisfied --- in the bilateral-commerce domain
and only there --- by a different sword: the bond the parties
have themselves posted, released or forfeit by an algorithm whose
execution they have themselves authorized. The Hobbesian
challenge is met on its own terms in the domain where the
substitution applies, and is not addressed in the domain where
the substitution does not apply. We do not argue that the
substitution generalizes outside bilateral commerce.

Schmitt's challenge is sharper and more interesting, and we
engage it directly rather than dismiss it. \citet{schmitt1922politicaltheology}
argues that any rule-system requires an extralegal decision on
the exception, and that the sovereign is whoever decides whether
a given case is or is not a normal application of the rules. The
bonded primitive, taken as a rule-system, would on Schmitt's
account require some sovereign to decide when it does not apply.
Our response is that the primitive is not a rule-system in
Schmitt's sense; it is a settlement primitive that runs in the
domain where its rules apply and is silent in the domain where
they do not. The Schmittian decision on the exception ---
whether a particular bonded commitment was procured under
duress, whether a particular outcome violates public policy,
whether a particular party lacks capacity --- is a decision the
primitive does not make and does not pretend to make. It is
made, where it is made at all, by external legal forums applying
their own constitutional and doctrinal apparatus to the bonded
commitment as evidentiary input (the companion evidence-and-law
paper develops this in detail). The exception is not absorbed by
the primitive; it is left where it has always been, in the
forums whose authority to decide on it derives from elsewhere.

The Schmittian challenge therefore lands, but lands at a
different level than its formulators may have intended. It is
correct that the primitive cannot decide its own exceptions;
that is precisely why the exception remains the province of
external forums, and why this paper is careful to bound its
substitution claim to bilateral commerce in normal application.
The non-discretion of the primitive within its domain, and its
silence outside its domain, are two faces of the same design
property: the primitive has no authority to extend itself, and
no authority to retract itself either. That property is what
makes coexistence with conventional sovereignty architectural
rather than negotiated.

\paragraph{Not a constitutional proposal.}
The paper does not propose constitutional reforms, alternative
governance structures, replacement of legislative or judicial
institutions, or any other prescriptive change to existing
political arrangements. The substrate change occurs at a layer
beneath constitutional choice, in the same way that the
electrical substrate underlying contemporary economies occurs at
a layer beneath constitutional choice. The constitutional
question of how to organize collective coercion in the remaining
domain is unaffected by the substrate change at that layer; what
is affected is the domain over which the question must be
answered.

\paragraph{Not a prediction.}
Finally, the paper does not predict that the state will wither,
that bilateral commerce will fully migrate to the bonded
substrate, or that the political consequences of such migration
will manifest in any particular timeframe or distribution. The
argument is conditional: under the supposition that the bonded
mechanism substitutes for sovereign-backed coercion in its
operable domain, the political-philosophical analysis of
coercion relocates in the manner described. Whether the
supposition obtains in practice is the empirical question of
adoption, which the present paper does not address.

% ════════════════════════════════════════════════════════════════════
\section{Conclusion}
\label{sec:conclusion}

This paper has argued that coercion is a substrate variable
analogous to subordination and coordination rent: structurally
load-bearing in the dominant tradition's account of bilateral
commercial enforcement, and substituted by pre-commitment when
the bonded primitive is the operable coordination architecture.
The substitution holds within a bounded domain (priceable,
bilateral, voluntary, schematizable obligations) and does not
extend beyond it (criminal harm, externality, public goods,
status, rights, defense). Within the domain, the
political-philosophical question of legitimate coercion becomes
irrelevant rather than corrected; outside the domain, the
question retains its full urgency, now operating over a bounded
rather than an unbounded scope.

The substantive observation is that the philosophical analysis
of coercion that has organized political thought from Hobbes
through Weber, Hart, and the contemporary literature has been
implicitly ranging over the entirety of social cooperation
because no architecture admitted bilateral commercial
enforcement without coercive backstop. With such an architecture
available, the analysis becomes bounded. This is not a refutation
of the prior tradition --- the analysis remains correct on its
own terms in the remaining domain --- but a relocation of where
the question of legitimate coercion must be answered. The state
retains its load-bearing functions; the question of how it
should exercise them retains its political-philosophical
urgency; what changes is the perimeter within which the
question operates.

We close with two observations of method, mirroring the closings
of the companion political-economy papers in the corpus. First,
the kernel is ideologically agnostic; the graph is the politics.
What the present paper adds to that formulation is that the
political-philosophical analysis of coercion, like the
political-economic analysis of subordination, becomes a
graph-tier question over a bounded substrate domain rather than
a substrate-tier inevitability over an unbounded one. The
politics is at the graph in a stronger sense: it is freed, in
the bounded domain, from a structural premise (coercive backstop
required) that prior architectures forced upon it. Second, the
relocation described here is recursive. Future substrates may
relocate further variables from architecture to choice, and the
boundary between bonded substitution and sovereign retention may
itself be re-examined as schema design extends what counts as
priceable. There is no terminal architecture; there is only the
next architecture whose variables we have not yet learned to see
as variables.

\section*{Acknowledgments}
The argument of this paper began as an intuition --- that the
mechanism's effectiveness varied inversely with the availability
of coercive alternatives, and that the dominant tradition's
analysis of law as text-plus-coercion required revision in the
substituting domain --- voiced in conversation. The intuition
was correct in direction but imprecise in mechanism, and this
paper is the unpacking of the more precise structural observation
to which it points. I thank the interlocutor whose intuition
prompted the unpacking, and the readers of earlier drafts who
pressed me to scope-bound the argument carefully and to resist
the triumphalist framings that the topic invites.


% ════════════════════════════════════════════════════════════════════
% References
% ════════════════════════════════════════════════════════════════════

\begin{thebibliography}{99}

\bibitem[Austin(1832)]{austin1832}
Austin, J.
\newblock \emph{The Province of Jurisprudence Determined}.
\newblock John Murray, London, 1832.

\bibitem[Elster(1979)]{elster1979ulysses}
Elster, J.
\newblock \emph{Ulysses and the Sirens: Studies in Rationality
  and Irrationality}.
\newblock Cambridge University Press, Cambridge, 1979.

\bibitem[Hart(1961)]{hart1961}
Hart, H.~L.~A.
\newblock \emph{The Concept of Law}.
\newblock Oxford University Press, Oxford, 1961.

\bibitem[Hobbes(1651)]{hobbes1651}
Hobbes, T.
\newblock \emph{Leviathan, or The Matter, Forme, and Power of a
  Common-Wealth Ecclesiasticall and Civill}.
\newblock Andrew Crooke, London, 1651.

\bibitem[Nozick(1969)]{nozick1969coercion}
Nozick, R.
\newblock Coercion.
\newblock In S.~Morgenbesser, P.~Suppes, and M.~White, editors,
  \emph{Philosophy, Science, and Method: Essays in Honor of Ernest
  Nagel}, pages 440--472. St.~Martin's Press, New York, 1969.

\bibitem[Nozick(1974)]{nozick1974anarchy}
Nozick, R.
\newblock \emph{Anarchy, State, and Utopia}.
\newblock Basic Books, New York, 1974.

\bibitem[Olson(1965)]{olson1965logic}
Olson, M.
\newblock \emph{The Logic of Collective Action: Public Goods and
  the Theory of Groups}.
\newblock Harvard University Press, Cambridge, MA, 1965.

\bibitem[Schmitt(1922)]{schmitt1922politicaltheology}
Schmitt, C.
\newblock \emph{Politische Theologie: Vier Kapitel zur Lehre von
  der Souver\"anit\"at}.
\newblock Duncker \& Humblot, Munich, 1922.
\newblock English: \emph{Political Theology: Four Chapters on the
  Concept of Sovereignty}, trans. G.~Schwab, MIT Press, 1985.

\bibitem[Srinivasan(2022)]{srinivasan2022network}
Srinivasan, B.
\newblock \emph{The Network State: How to Start a New Country}.
\newblock Self-published, 2022.

\bibitem[Stigler(1971)]{stigler1971economic}
Stigler, G.~J.
\newblock The Theory of Economic Regulation.
\newblock \emph{Bell Journal of Economics and Management Science},
  2(1):3--21, 1971.

\bibitem[Weber(1919)]{weber1919}
Weber, M.
\newblock \emph{Politik als Beruf}.
\newblock Duncker \& Humblot, Munich, 1919.
\newblock English: ``Politics as a Vocation'', in H.~H.~Gerth and
  C.~Wright Mills, eds., \emph{From Max Weber: Essays in
  Sociology}, Oxford University Press, 1946.

\bibitem[Weber(1922)]{weber1922economy}
Weber, M.
\newblock \emph{Wirtschaft und Gesellschaft}.
\newblock J.~C.~B.~Mohr, T\"ubingen, 1922.
\newblock English: \emph{Economy and Society: An Outline of
  Interpretive Sociology}, G.~Roth and C.~Wittich, eds.,
  University of California Press, Berkeley, 1978.

\bibitem[Pettit(1997)]{pettit1997republicanism}
Pettit, P.
\newblock \emph{Republicanism: A Theory of Freedom and Government}.
\newblock Oxford University Press, Oxford, 1997.

\bibitem[Pettit(2012)]{pettit2012on}
Pettit, P.
\newblock \emph{On the People's Terms: A Republican Theory and
  Model of Democracy}.
\newblock Cambridge University Press, Cambridge, 2012.

\bibitem[Friedman(1973)]{friedman1973machinery}
Friedman, D.~D.
\newblock \emph{The Machinery of Freedom: Guide to a Radical
  Capitalism}.
\newblock Harper \& Row, New York, 1973.

\bibitem[Williamson(1985)]{williamson1985}
Williamson, O.~E.
\newblock \emph{The Economic Institutions of Capitalism}.
\newblock Free Press, New York, 1985.

\end{thebibliography}

\end{document}
